% TOP_PROBLEM: {"id": 30007015, "title": "Bounded-metric anticanonical nonvanishing for klt pairs", "category_id": 5, "category": {"id": 5, "name": "algebraic_geometry", "display_name": "Algebraic Geometry", "description": "Geometric objects defined by polynomial equations.", "slug": "algebraic-geometry", "order_index": 5, "created_at": "2026-07-31T15:26:25.671Z"}, "status": "open"}
% FIRST_POSTED: 2026-10-06
% ATTEMPT_STATUS: solved
% OPENAI_CLAIM: {"schema_version":1,"source_repo":"https://github.com/openai/math","source_commit":"adc7f1241b42e322a6451854ab7e4b4c146bf78a","release_date":"2026-10-06","match":"full","resolution":"proved","independently_reviewed":false,"second_pass":"agreed","families":[{"family":"068","title":"Anticanonical nonvanishing in every dimension","subject":"Algebraic and complex geometry","manuscripts":[{"title":"Bounded anticanonical metrics on klt pairs and torus quotients","pdf_path":"preprints/Bounded-anticanonical-metrics-on-klt-pairs-and-torus-quotients-September-26-2026/main.pdf","readme_path":"preprints/Bounded-anticanonical-metrics-on-klt-pairs-and-torus-quotients-September-26-2026/README.md","date":"2026-09-26","theorem_ref":"thm:klt (source/sections/introduction.tex:15)","statement_sources":[{"path":"preprints/Bounded-anticanonical-metrics-on-klt-pairs-and-torus-quotients-September-26-2026/build/source/sections/introduction.tex","line":15,"label":"thm:klt"}]}],"lean":null,"reasoning_trace":null}]}
\documentclass[11pt]{article}
\usepackage[margin=25mm]{geometry}
\usepackage{fontspec}
\setmainfont{Latin Modern Roman}
\usepackage{amsmath,amssymb,mathtools}
\usepackage{unicode-math}
\setmathfont{Latin Modern Math}
\usepackage{enumitem}
\usepackage{xurl,hyperref,bookmark}
\hypersetup{colorlinks=true,urlcolor=blue}
\setcounter{secnumdepth}{0}
\setlength{\parindent}{0pt}
\setlength{\parskip}{5pt}
\setlength{\emergencystretch}{3em}
\begin{document}
\section*{Bounded-metric anticanonical nonvanishing for klt pairs}

\subsection{OpenAI's claim}
A nef anti-log-canonical divisor with bounded semipositive metric is claimed to acquire a section when a resolution has nonzero structure-sheaf Euler characteristic; the torus-quotient application gives naturally invariant sections. The cited manuscript argues for that result; its relation to this problem is specified below.

\subsection{Mathematical statement of the claimed result}
\paragraph{Mathematical claim 1.} For a projective rational-factorial klt complex pair $(W,D)$ with effective rational boundary, set $P=-(K_W+D)$ and assume $P$ nef. Require a smooth projective resolution $\rho:\widetilde W\to W$ on which $\rho^*P$ has a semipositive singular Hermitian metric with locally bounded weights, and
\[
\chi(\widetilde W,\mathcal O_{\widetilde W})\ne0.
\]
OpenAI claims a positive Cartier multiple with $H^0(W,\mathcal O_W(mP))\ne0$. Locally bounded weights need not be continuous. Nefness alone is not the stated hypothesis; the bounded metric and Euler-characteristic condition are retained.

\subsection{Outline of the argument}
Manuscript 1 (family 068, 2026-09-26): The nonzero Euler characteristic first supplies twisted differentials, and cotangent-subsheaf and birational arguments convert them into the specified line. A stable projective torus quotient is built with the same bounded-metric hypotheses, so its section pulls back as an invariant anticanonical plurisection.

\subsection{Scope relative to this problem}
Manuscript 1 (family 068, 2026-09-26): The problem selects the mathematical statement summarized here, including its stated hypotheses and parameter range. The cited theorem gives a claimed proof of that assertion in full.

\subsection{Formal verification}
Family 068 has no Lean scope page in this snapshot; no formalization of these manuscript results is listed.

This site has not executed the Lean files or checked their dependencies. These descriptions report the scope of the linked declarations, rather than a new verification of the proofs.

\subsection{Status caveat}
OpenAI attributes these results to an internal model. Its release notice cautions that results lacking full formalization can contain mistakes. This site has not independently assessed the proofs; a solved label records the scope of the reported claim.

\subsection{Sources}
\begin{itemize}
\item OpenAI, \emph{Bounded anticanonical metrics on klt pairs and torus quotients}, 2026-09-26; citation key \texttt{OAI:Bounded-anticanonical-metrics-on-klt-pairs-and-torus-quotients-September-26-2026}. \href{https://github.com/openai/math/blob/adc7f1241b42e322a6451854ab7e4b4c146bf78a/preprints/Bounded-anticanonical-metrics-on-klt-pairs-and-torus-quotients-September-26-2026/main.pdf}{PDF (GitHub)}; \href{https://github.com/openai/math/raw/adc7f1241b42e322a6451854ab7e4b4c146bf78a/preprints/Bounded-anticanonical-metrics-on-klt-pairs-and-torus-quotients-September-26-2026/main.pdf}{PDF (direct)}; \href{https://github.com/openai/math/blob/adc7f1241b42e322a6451854ab7e4b4c146bf78a/preprints/Bounded-anticanonical-metrics-on-klt-pairs-and-torus-quotients-September-26-2026/README.md}{Manuscript page and citation}; \href{https://github.com/openai/math/blob/main/preprints/Bounded-anticanonical-metrics-on-klt-pairs-and-torus-quotients-September-26-2026/main.pdf}{Latest version}; Statement source: \href{https://github.com/openai/math/blob/adc7f1241b42e322a6451854ab7e4b4c146bf78a/preprints/Bounded-anticanonical-metrics-on-klt-pairs-and-torus-quotients-September-26-2026/build/source/sections/introduction.tex\#L15}{thm:klt}.
\item \href{https://github.com/openai/math/blob/adc7f1241b42e322a6451854ab7e4b4c146bf78a/CONTENTS.md}{OpenAI catalogue entry}.
\end{itemize}
\end{document}
